
\usepackage{titlesec}
\usepackage{picture}
% ---- 行の高さに自動追従する左肩の細四角 ----
\newcommand{\TitleSquareToText}[2][1.05ex]{%
  \begingroup
    \leavevmode
    \setbox0=\hbox{\strut}%
    \textcolor{#2}{\rule[-\dp\strutbox]{#1}{\dimexpr\ht\strutbox+\dp\strutbox\relax}}%
  \endgroup
}

% ---------------- Part（下線=フル幅） ----------------
\titleformat{\part}[display]
  {\normalfont\LARGE\bfseries\centering\color{appleInk}}
  {\small\scshape\color{appleGray}第\thepart 章}
  {0.3ex}
  {}%
  [{%
     \vspace{0.6ex}%
     {\color{appleLine}\titlerule[1.2pt]}%
     \nobreak\vspace{0.25ex}%
     {\color{appleLine}\titlerule[0.6pt]}%
  }]
  \begin{comment}
% ---------------- Section（下線=66%） ----------------
\titleformat{\section}[block]
  {\normalfont\Large\bfseries\color{appleInk}}
  {}{0pt}
  {%
    \TitleSquareToText[1.05ex]{SectionTag}\hspace{0.6em}%
    \thesection\quad
  }%
  [{%
     \vspace*{-0.48\baselineskip}% ← 線を上に引き上げて詰める
     \noindent{\color{appleLine}\rule{0.66\linewidth}{0.8pt}}%
     \nobreak\vspace{-0.10\baselineskip}% ← 線の下側（本文との距離）
  }]
\end{comment}
% ---------------- Subsection（下線=50%） ----------------
\titleformat{\section}[block]
  {\normalfont\large\bfseries\color{appleInk}}
  {}{0pt}
  {%
    \TitleSquareToText[0.95ex]{SubsectionTag}\hspace{0.55em}%
    \thesection\quad
  }%
  [{%
     \vspace*{-0.52\baselineskip}% ← もう少しだけ控えめに詰める
     \noindent{\color{appleLine}\rule{0.50\linewidth}{0.6pt}}%
     \nobreak\vspace{-0.18\baselineskip}%
  }]

% ---------------- Subsubsection（下線=42%） ----------------
\titleformat{\subsection}[block]
  {\normalfont\large\bfseries\color{appleInk}}
  {}{0pt}
  {%
    \TitleSquareToText[0.90ex]{SubsubsectionTag}\hspace{0.50em}%
    \thesubsection\quad
  }%
  [{%
     \vspace*{-0.40\baselineskip}% ← さらに控えめ
     \noindent{\color{appleLine}\rule{0.42\linewidth}{0.5pt}}%
     \nobreak\vspace{-0.16\baselineskip}%
  }]

% =========================
% 柱・ページ番号
% =========================
\usepackage{fancyhdr}

\newcommand{\CurrentPartTitle}{}
\newcommand{\CurrentSectionTitle}{}

\makeatletter
\let\origpart\part
\renewcommand{\part}{%
  \@ifstar{\part@star}{\@ifnextchar[{\part@opt}{\part@noopt}}%
}

\def\part@noopt#1{%
  \origpart{#1}%
  \gdef\CurrentPartTitle{第\thepart 章\quad #1}%
}

\def\part@opt[#1]#2{%
  \origpart[#1]{#2}%
  \gdef\CurrentPartTitle{第\thepart 章\quad #1}%
}

\def\part@star#1{%
  \origpart*{#1}%
  \gdef\CurrentPartTitle{#1}%
}

\let\origsection\section
\renewcommand{\section}{%
  \@ifstar{\section@star}{\@ifnextchar[{\section@opt}{\section@noopt}}%
}

\def\section@noopt#1{%
  \origsection{#1}%
  \gdef\CurrentSectionTitle{\thesection\quad #1}%
}
\def\section@opt[#1]#2{%
  \origsection[#1]{#2}%
  \gdef\CurrentSectionTitle{\thesection\quad #1}%
}
\def\section@star#1{%
  \origsection*{#1}%
  \gdef\CurrentSectionTitle{#1}%
}
\makeatother

\pagestyle{fancy}
\fancyhf{}
\renewcommand{\headrulewidth}{0pt}

\fancyhead[LE]{%
  \footnotesize
  p.\,\thepage\hspace{1.2\zw}\CurrentPartTitle
}

\fancyhead[RO]{%
  \footnotesize
  \CurrentSectionTitle\hspace{1\zw}|\hspace{1\zw}p.\,\thepage
}